\section{Reproducibility, source provenance, and claim ledger}

\subsection{What is in the archive}
The archive contains this PDF and its complete LaTeX source, bibliography,
figure-generation source and rendered figures, a native integration patch,
the proposed native files, the standalone forest solvers and their separate
verifier, focused tests, recorded audit and benchmark output, positive and
negative source packets, and integration/reproduction instructions.
\code{README.md} is the entry point; \code{REPRODUCE.md} gives exact
commands.  \code{code_manifest.json} records the proposed code identities,
and \code{SHA256SUMS.txt} records archive file digests.

The native patch is relative to
\begin{center}\small\basecommit.\end{center}
It is not relative to whichever commit the remote branch later names
\code{main}.  The package's patch validation reconstructs the pinned
original files, applies the patch, and compares the result byte-for-byte
with the supplied proposed files.  A fresh complete clone is preferable
for the full regression suite because older tests import exact reference
implementations and fixtures from historical \code{reports} directories.
Those fixture identities are listed in \code{historical_fixtures.json}
under \code{results/package_validation/}.

The standalone finite tools require only Python's standard library.
The native tests additionally require the pinned ProveIt source tree;
Regina is needed for the comparison audit, not for the replay of the seven
saved positive packets or three source-authenticated negative packets.
Matplotlib is needed only to regenerate the figures.  The article uses
standard pdfLaTeX packages and BibTeX.  Exact commands and recorded versions
are included rather than an implicit environment assumption.

\subsection{Evidence and theorem mapping}
\begin{longtable}{>{\raggedright\arraybackslash}p{0.29\textwidth}>{\raggedright\arraybackslash}p{0.61\textwidth}}
\caption{How each substantive claim is supported.}\label{tab:claims}\\
\toprule
Claim & Proof or independent evidence\\
\midrule\endfirsthead
\toprule Claim & Proof or independent evidence\\\midrule\endhead
Native graph root reachability & Lemma~\ref{lem:rooted}: connected source
one-skeleton with positive edges replaced by root stubs.\\
Geometric reduction of all top handles & Lemma~\ref{lem:dual-cancel} and
Theorem~\ref{thm:native}; actual belt occurrences, not signed coefficients.\\
Retention of native exposed ports & Lemma~\ref{lem:ports}; a relative
boundary-ball argument for authenticated exposed patches.\\
Polynomial encoded graph stage & Bounds on native chart records and
arithmetic in Section~\ref{sec:geometry}; no normal-piece expansion.\\
Stopping characterization & Prior incoming stopping theorem, re-proved
for the stated raw update in Section~\ref{sec:forests}.\\
Complete constrained forest search & Theorem~\ref{thm:search}; independent
branch-covering proof-tree replay.\\
Polynomial sparse-column selection & Reachability and source-SCC lemmas,
Rado's criterion, and a proved shortest-path augmentation.\\
Native total-column bound and parity & Propositions~\ref{prop:three}
and~\ref{prop:cycle}; triangular face regions and the handle-chain identity.\\
85 to 92 complete native cuts & Pinned baseline/native audit, preserved
censuses, 92 independent replays, seven Regina comparisons.\\
Three interface kernels admit no tree & Saved authenticated negative
packets and independent reconstruction of their nondeletable stopping rows.\\
Finite algorithm agreement & 26,995 exhaustive and 7,000 seeded random
oracle comparisons; separate positive and negative checkers.\\
Finite enumeration speed comparison & Specified $2^k$-tree family,
complete-query timings, raw paired rounds and structural counters.\\
Native timing crossover & Six supplied-cut sizes and five randomized
paired rounds per size; scope and regressions retained.\\
General quasi-polynomial unknot recognition & Not proved here.  Discovery,
global-size, pattern-transport, and terminal-completeness obligations remain.\\
\bottomrule
\end{longtable}

\subsection{A small inspectable finite certificate}
The following example is the finite form of the two-choice construction:
\begin{verbatim}
{
  "vertices": 2,
  "columns": 1,
  "edges": [[0, 1], [0, 1]],
  "rows": [[[0, 1]], [[0, 3]]]
}
\end{verbatim}
The tree consists of row 1.  After deleting it, row 0 is removed on
column 0.  A positive finite certificate therefore records the tree
$[1]$ and the peel pair $[0,0]$.  Changing the selected tree to $[0]$
without changing the source fails verification: the surviving count is
three, not one.  If both counts are three, the two stopping singleton
groups have union of graphic rank one, smaller than their group count
two.  That rank defect certifies \code{NO_FOREST}; it does not certify
that the associated manifold is not a ball.

The delivered demonstration generates several such positive and negative
certificates, replays them with the independent checker, and rejects
mutated variants.  The full source packets are larger because they also
authenticate the triangulation, coordinates, geometry, and interfaces.

\subsection{Input validation versus geometric interpretation}
A graph/matrix result can be correct as a finite result even when no
three-manifold realizes the input.  Conversely, a source-authenticated
topological result needs the entire geometric chain of implications.
The manuscript separates these levels to make review practical:
\begin{enumerate}
\item The native source validator establishes the manifold and normal-vector
  contract and authenticates the complementary charts.
\item The geometric dual theorem identifies valid top cancellations and
  the retained native port data.
\item The finite matrix theorem decides exactly the stated raw
  forest-completion problem on those authenticated occurrences.
\item The boundary kernel certifies a genus census when its additional
  attachment criteria succeed.
\end{enumerate}
The last level can remain inconclusive even when the first three succeed.
The ten-core audit demonstrates this distinction directly.  A negative
finite certificate must not cross these levels as a knot verdict.

There is a second distinction between actual counts and saturated counts.
Capping at two is exact for every raw pivot and stopping predicate in
this paper.  It is not exact for parity: an actual count of three is odd
but becomes two.  Proposition~\ref{prop:cycle} uses authenticated uncapped
counts.  The source parity audit follows that convention separately from
the finite solver's capped matrix representation.

\subsection{Attribution and limits of the research claim}
The topological cancellation principle is classical.  Rado's independent
representatives theorem and Edmonds' matroid-intersection framework are
classical.  Compressed normal-surface cutting is established in the
Agol--Hass--Thurston and Lackenby literature.  The stopping-set criterion
and its raw-cancellation extension already occur in the inspected incoming
research.  The contribution claimed here is their specific connection to
the native source geometry, the retained interface proof, the exact
constrained tree-selection formulation and sparse-column solution, the
native column constraints, and the tested certificate implementation.

We did not perform an exhaustive historical priority search for every
combinatorial lemma, and these results have not undergone external peer
review or proof-assistant verification.  The claims are stated so that
they can be checked or refuted independently.  The archive supplies the
counterexamples, exact inputs, patch, raw timings, and negative evidence
needed for that review.  Numerical experiments support implementation
correctness and practical scope; they do not replace the proofs or imply
the unresolved global asymptotic bound.
